\section{Review of Projectile Motion}
\subsection{basic}
A simple projectile is an object that has a single, \textit{non-uniform} force acting on it. This force must be a 
\textit{non-contact} force between objects that are not in contact. Examples include gravitational forces, magnetic forces,
and electric forces. \\

Some common question types are:
\begin{enumerate}
    \item An object dropped
    \item A soccer ball is kicked
    \item A bullet is fired from a gun 
    \item An electron is moving through a uniform electric field (gravity is negligible)
\end{enumerate}

Remember: the object accelerates only in the direction of the net force. In any direction perpendicular to the net force, 
the object maintains a constant velocity. 

Some important points to remember:
\begin{itemize}
    \item At maximum height, all projectiles have a \textbf{vertical} velocity equal to \textbf{zero}.
    \item When an object starts and ends at the same vertical location, then $\vec{\Delta d_{y}} = 0$.
    \item When an object is dropped or launched horizontally, then $\vec{v_{y1}} = 0$.
\end{itemize}

\subsection{Special formula}
\begin{equation} \label{eq:Special}
    R = \frac{ v_{i}^2 * sin{2\theta} }{g} 
\end{equation}
\begin{center}
    $R$ is the range of the projectile (horizontal distance in m) \\
    $v_{i}$ is the launch speed of the projectile in $(m/s)$ \\
    $\theta$ is the launch angle of the projectile \\
    $g$ is the acceleration due to gravity \\
\end{center}

The formula \ref{eq:Special} can only be used if the projectile \textcolor{red}{starts and ends at the same vertical location}.
\newpage
\subsection{An example question}
Remember, you should always label your direction conventions.

\mypic{pictures/projectileReview.png}{}{0.8}
